\section{Phương pháp nghiên cứu}

Đồ án kết hợp nghiên cứu tài liệu, thực nghiệm mô hình và xây dựng hệ thống. Các
tài liệu về Edge AI, YOLO, cắt tỉa, chưng cất, lượng tử hóa, TensorRT, hệ thống
phân tán, AI Agent và Domain-Driven Design được tổng hợp để hình thành cơ sở lý
thuyết.

Đối với mô hình, phương pháp thực nghiệm có kiểm soát được sử dụng: giữ cố định
tập dữ liệu, cách chia tập, kích thước ảnh và tiêu chí đo; sau đó thay đổi từng
yếu tố như phương án phục hồi, trọng số chưng cất hay precision. Thực nghiệm sơ
bộ trên tập nhỏ được chạy trước để kiểm chứng pipeline, lặp với nhiều seed để
đánh giá độ biến thiên, trước khi chạy tốn kém trên tập mục tiêu.

Đối với nền tảng, đồ án áp dụng Domain-Driven Design và kiến trúc port--adapter,
hiện thực theo mô hình Cloud Control Plane -- Edge Agent. Tính đúng đắn được kiểm
chứng bằng kiểm thử tự động chạy mã sản phẩm của Cloud và Agent với hạ tầng giả
lập trong bộ nhớ, bao gồm các kịch bản triển khai thành công, xung đột, lỗi pull
image, container bị crash, telemetry và xác thực; sau đó được kiểm chứng trên hệ
thống triển khai thật.
